\problem{4.35}{Spin measurement probabilities during Larmor precession}
Example 4.3 gives, for real initial coefficients,
$\chi(t)=(\cos(\alpha/2)\ee^{\ii\gamma B_0t/2},
\sin(\alpha/2)\ee^{-\ii\gamma B_0t/2})^T$.
The spin expectations are
\begin{equation}
\avg{S_x}=\frac\sh2\sin\alpha\cos\omega_Lt,\quad
\avg{S_y}=-\frac\sh2\sin\alpha\sin\omega_Lt,\quad
\avg{S_z}=\frac\sh2\cos\alpha,\quad \omega_L=\gamma B_0.
\end{equation}
For spin one-half,
$P(S_i=+\sh/2)=\frac12+\avg{S_i}/\sh$.
Therefore
\begin{align}
\boxed{P_x(+)}&=\frac12[1+\sin\alpha\cos(\omega_Lt)],\\
\boxed{P_y(+)}&=\frac12[1-\sin\alpha\sin(\omega_Lt)],\\
\boxed{P_z(+)}&=\frac12(1+\cos\alpha)=\cos^2(\alpha/2).
\end{align}

\problem{4.36}{Exact spin dynamics in an oscillating longitudinal field}
For $\mathbf B=B_0\cos(\omega t)\hat z$,
$H(t)=-\gamma B_0\cos(\omega t)S_z$,
so
\begin{equation}
\boxed{H(t)=-\frac{\gamma\sh B_0}{2}
\cos(\omega t)\begin{pmatrix}1&0\\0&-1\end{pmatrix}}.
\end{equation}
The Hamiltonian commutes with itself at all times,
so the evolution is an ordinary diagonal phase:
\begin{equation}
\boxed{U(t)=
\begin{pmatrix}\ee^{\ii\vartheta(t)}&0\\
0&\ee^{-\ii\vartheta(t)}\end{pmatrix}},
\quad
\vartheta(t)=\frac{\gamma B_0}{2\omega}\sin(\omega t).
\end{equation}
With initial $\chi_{x,+}=(1,1)^T/\sqrt2$,
$\chi(t)=
(\ee^{\ii\vartheta},\ee^{-\ii\vartheta})^T/\sqrt2$.
Projection onto $\chi_{x,-}=(1,-1)^T/\sqrt2$
gives amplitude $\ii\sin\vartheta$ and
\begin{equation}
\boxed{P(S_x=-\sh/2,t)=
\sin^2\left[\frac{\gamma B_0}{2\omega}\sin(\omega t)\right]}.
\end{equation}
Complete inversion requires the maximal accessible
$|\vartheta|\ge\pi/2$, so
$\boxed{|B_0|_{\min}=\pi|\omega|/|\gamma|}$,
for nonzero oscillation frequency.

\problem{4.37}{Confirming triplets and the singlet using ladder operators}
Write $\ket{1,0}=(\ket{\uparrow\downarrow}+
\ket{\downarrow\uparrow})/\sqrt2$.
With $S_-=S_-^{(1)}+S_-^{(2)}$,
\begin{equation}
S_-\ket{1,0}=
\frac1{\sqrt2}(\sh\ket{\downarrow\downarrow}
+\sh\ket{\downarrow\downarrow})
=\boxed{\sh\sqrt2\,\ket{1,-1}}.
\end{equation}
For the singlet
$\ket{0,0}=(\ket{\uparrow\downarrow}-
\ket{\downarrow\uparrow})/\sqrt2$,
the terms cancel,
$\boxed{S_\pm\ket{0,0}=0}$.
Since $S_z\ket{0,0}=0$ as well,
$S^2=S_+S_-+S_z^2-\sh S_z$
gives $\boxed{S^2\ket{0,0}=0}$.
For a triplet state the same expression gives
$S^2=2\sh^2$.

\problem{4.38}{Allowed baryon and meson spins}
For three spin-one-half constituents,
first couple two spins to $s_{12}=0$ or $1$.
Then $0\otimes\frac12=\frac12$,
while $1\otimes\frac12=\frac32\oplus\frac12$.
Thus
\begin{equation}
\boxed{\text{three quarks: }s=\frac12
\text{ or }\frac32}.
\end{equation}
For a quark and antiquark with spin one-half,
$\frac12\otimes\frac12=0\oplus1$:
\begin{equation}
\boxed{\text{two constituents: }s=0\text{ or }1}.
\end{equation}
These results assume zero orbital angular momentum.
Additional color, flavor and exchange constraints
are needed to classify physical hadrons.

\problem{4.39}{Clebsch--Gordan construction for two spin halves}
For $s_1=s_2=1/2$, the coupled triplet is
\begin{align}
\ket{1,1}&=\ket{\uparrow\uparrow},\\
\ket{1,0}&=\frac{\ket{\uparrow\downarrow}+
\ket{\downarrow\uparrow}}{\sqrt2},\\
\ket{1,-1}&=\ket{\downarrow\downarrow},
\end{align}
and the singlet is
$\ket{0,0}=(\ket{\uparrow\downarrow}-
\ket{\downarrow\uparrow})/\sqrt2$.
The corresponding Clebsch--Gordan coefficients
are $1,\;1/\sqrt2,\;1/\sqrt2,\;1$
for the triplet, and $+1/\sqrt2,-1/\sqrt2$
for the singlet. Each expansion is normalized.
Their inner products vanish across the
triplet/singlet sectors, confirming the
decomposition $2\otimes2=3\oplus1$.

\problem{4.40}{Probabilities for adding unequal spins}
(a) For $s_1=1,s_2=2$, the state $\ket{3,1}$
contains contributions with
$m_1+m_2=1$:
\begin{equation}
\ket{3,1}=
\sqrt{\frac25}\ket{1,1;2,0}
+\sqrt{\frac8{15}}\ket{1,0;2,1}
+\sqrt{\frac1{15}}\ket{1,-1;2,2}.
\end{equation}
Hence measuring $S_z^{(2)}$ yields
\begin{equation}
\boxed{m_2=0,1,2\quad\text{with probabilities }
2/5,\ 8/15,\ 1/15}.
\end{equation}
(b) A hydrogen electron in $\psi_{510}$
has orbital $\ell=1,m_\ell=0$; with spin
down $m_s=-1/2$. The possible total
$j$ values are $3/2$ and $1/2$.
The relevant Clebsch--Gordan coefficients
are $\sqrt{2/3}$ and $\sqrt{1/3}$,
respectively. Thus
\begin{equation}
\boxed{P(J^2=\tfrac{15}{4}\sh^2)=\frac23,\qquad
P(J^2=\tfrac34\sh^2)=\frac13}.
\end{equation}

\problem{4.41}{Total spin commutes with its $z$ component}
With $\mathbf S=\mathbf S_1+\mathbf S_2$,
$S^2=S_1^2+S_2^2+2\mathbf S_1\cdot\mathbf S_2$.
Each individual $S_a^2$ commutes with all
components of that particle's spin.
For the cross term,
\begin{equation}
[S_1\cdot S_2,S_{1k}]
=\sum_i[S_{1i},S_{1k}]S_{2i}
=\ii\sh\sum_{i,j}\epsilon_{ikj}S_{1j}S_{2i}.
\end{equation}
Rearranging indices yields the vector result
\begin{equation}
\boxed{[S^2,\mathbf S_1]
=2\ii\sh(\mathbf S_1\times\mathbf S_2)},
\quad
\boxed{[S^2,\mathbf S_2]
=-2\ii\sh(\mathbf S_1\times\mathbf S_2)}.
\end{equation}
Adding them gives $[S^2,\mathbf S]=0$,
and therefore $\boxed{[S^2,S_z]=0}$.
The corresponding commutator with an
\emph{individual} $S_{1z}$ is generally nonzero.

\problem{4.42}{Minimal coupling and the quantum Lorentz-force law}
Let $\boldsymbol\pi=\mathbf p-q\mathbf A$,
$\mathbf v=\boldsymbol\pi/m$ and
$H=\boldsymbol\pi^2/(2m)+q\phi$.
Using $[r_i,\pi_j]=\ii\sh\delta_{ij}$,
\begin{equation}
\boxed{\frac{\dd\avg{\mathbf r}}{\dd t}
=\frac{\avg{\boldsymbol\pi}}m=\avg{\mathbf v}}.
\end{equation}
The kinetic commutator is
$[\pi_i,\pi_j]=\ii q\sh\,\epsilon_{ijk}B_k$.
Applying Ehrenfest to $\pi_i$ and
including explicit $\partial_t\mathbf A$,
one obtains the symmetric-ordered result
\begin{equation}
m\frac{\dd\avg{\mathbf v}}{\dd t}
=q\avg{\mathbf E}
+\frac{q}{2m}
\avg{\boldsymbol\pi\times\mathbf B-
\mathbf B\times\boldsymbol\pi}.
\end{equation}
Writing $\boldsymbol\pi=\mathbf p-q\mathbf A$
recovers the equivalent textbook formula with
$\frac q{2m}(\mathbf p\times\mathbf B-
\mathbf B\times\mathbf p)-(q^2/m)\mathbf A\times\mathbf B$.
For spatially uniform $\mathbf E,\mathbf B$,
operator-ordering issues disappear:
\begin{equation}
\boxed{m\,\frac{\dd\avg{\mathbf v}}{\dd t}
=q(\mathbf E+\avg{\mathbf v}\times\mathbf B)}.
\end{equation}

\problem{4.43}{Landau levels plus confinement along the field}
With $\mathbf A=(B_0/2)(-y,x,0)$
and $\phi=Kz^2$, the electromagnetic fields are
\begin{equation}
\boxed{\mathbf B=B_0\hat z,\qquad
\mathbf E=-2Kz\hat z}.
\end{equation}
The Hamiltonian separates into a transverse
Landau problem and a $z$ oscillator:
\begin{equation}
H_\perp=\frac{(p_x+qB_0y/2)^2+
(p_y-qB_0x/2)^2}{2m},\qquad
H_z=\frac{p_z^2}{2m}+qKz^2.
\end{equation}
The transverse kinetic momenta satisfy
$[\pi_x,\pi_y]=\ii q\sh B_0$, giving
$E_\perp=\sh\omega_c(n_1+1/2)$,
$\omega_c=|qB_0|/m$.
For $qK>0$,
$H_z$ is an oscillator with
$\omega_z=\sqrt{2qK/m}$, so
\begin{equation}
\boxed{E_{n_1,n_2}=
\sh\omega_c(n_1+\tfrac12)+
\sh\omega_z(n_2+\tfrac12)}.
\end{equation}
The Landau levels have extensive transverse
degeneracy on an unbounded plane.
If $qK<0$, longitudinal motion is not
confined and the quoted discrete ladder does
not apply.

\problem{4.44}{Gauge covariance of the Schrödinger equation}
For real $\Lambda(\mathbf r,t)$, define
$\mathbf A'=\mathbf A+\nabla\Lambda$,
$\phi'=\phi-\partial_t\Lambda$,
and $\Psi'=\ee^{\ii q\Lambda/\sh}\Psi$.
Act on the new state with the
gauge-transformed covariant momentum:
\begin{equation}
(-\ii\sh\nabla-q\mathbf A')\Psi'
=\ee^{\ii q\Lambda/\sh}
(-\ii\sh\nabla-q\mathbf A)\Psi.
\end{equation}
Applying the operator again proves the same
identity for the kinetic-energy square.
The time derivative introduces
\begin{equation}
\ii\sh\partial_t\Psi'
=\ee^{\ii q\Lambda/\sh}
(\ii\sh\partial_t\Psi-
q\,\partial_t\Lambda\,\Psi).
\end{equation}
The extra term is precisely canceled by
$q\phi'=q(\phi-\partial_t\Lambda)$:
\begin{equation}
\boxed{\ii\sh\partial_t\Psi'=H'\Psi'}.
\end{equation}
Hence predictions are gauge invariant
provided the wavefunction also transforms.

\problem{4.45}{Energy shifts on a magnetic-flux-threaded ring}
A charge $q$ is constrained to a ring of
radius $b$ encircling magnetic flux $\Phi$.
Outside the solenoid $\mathbf B=0$, but
the circulating vector potential is
$A_\phi=\Phi/(2\pi b)$.
Minimal coupling gives
\begin{equation}
\boxed{H=\frac1{2mb^2}
\left(-\ii\sh\frac{\dd}{\dd\phi}
-\frac{q\Phi}{2\pi}\right)^2}.
\end{equation}
Periodic eigenfunctions
$\psi_n(\phi)=\ee^{\ii n\phi}/\sqrt{2\pi}$
yield
\begin{equation}
\boxed{E_n(\Phi)=\frac{\sh^2}{2mb^2}
\left(n-\frac{q\Phi}{2\pi\sh}\right)^2},
\qquad n\in\mathbb Z.
\end{equation}
The spectrum is periodic under
$\Phi\mapsto\Phi+\Phi_0$,
$\Phi_0=h/|q|$, after relabeling $n$.
On two distinct paths around the solenoid,
the relative phase is
$\boxed{\Delta\varphi=q\Phi/\sh}$,
which can shift an interference pattern
even when the particle never visits a region
with nonzero magnetic field.
